% Copyright 2006, 2008, 2009 W.I.S.V. Christiaan Huygens

\documentclass[a4paper]{artikel3}

\usepackage[english]{babel}
\usepackage{epsfig}
\usepackage{fullpage}
\usepackage[T1]{fontenc}

\newcommand{\szp}{\ensuremath{\mathcal{SZP}}}
\newcommand{\cmd}[1]{\texttt{#1}}
\newcommand{\file}[1]{\texttt{#1}}
\newcommand{\key}[1]{\textbf{#1}}
\newcommand{\strong}[1]{\textbf{#1}}
\newcommand{\syntax}[1]{\textbf{Syntax}\\\cmd{#1}}
\newenvironment{citemize}
  {\begin{list}{$\bullet$}{\topsep 0cm \itemsep 0cm}}
  {\end{list}}

\parskip 0.2cm
\parindent 0cm

\title{ Sub Zero Programming 4.0 \\ Team Manual }

\begin{document}

\begin{titlepage}
\begin{center}

\vspace{2cm}
{\Huge Sub Zero Programming} \\
\vspace{1cm}
{\Huge Team Manual} \\
%\vspace{2cm}
%{\Large July 2004} \\
\vspace{5cm}

\epsfig{file=szplogo.eps}

\end{center}
\end{titlepage}

% \tableofcontents

\bigskip

\vspace{-0.8cm}
\section{Introduction}

TeamTool is your interface with the jury during the contest. It allows you
to submit solutions and send messages to the jury. It consists of two parts:
\cmd{submit}, a command line tool to submit solutions and TeamTool, a web
interface to get contest information and communicate with the jury.

\vspace{-0.8cm}
\section{\cmd{submit}}

\label{sec:submit}

When you solved a problem, you have to submit it, so it can be judged. You
can do this using TeamTool or by using the \cmd{submit} command. The latter
is recommended.

The \cmd{submit} program is a command-line tool. You should start it from a
command shell.

\syntax{submit <problem-letter> L<language-number> <sourcefile>}

You can start submit without any parameters to get a list of the problem
letters and language numbers to use.

When invoked properly, the program shows details about your command, and
asks you to confirm the submission. Respond by typing `Y' if you want to
continue.

The program will give you a warning if you specified a source file with a
non-standard file name extension for your programming language. It also
warns you if you specified a file that has not been modified in the last 5
minutes. The purpose of these warnings is to prevent you from accidentally
submitting a wrong file. If you choose to continue the submission, the
system will handle your submission normally.

\strong{Notes}
\begin{citemize}
\item Make sure to submit only the source code of your program; executable
files cannot be submitted. The jury will compile your source code to
generate an executable program.

\item It is not possible to use multiple source files in one submission; you
must put all the source code for your program in a single file.
\end{citemize}

\vspace{-0.8cm}
\section{TeamTool}

TeamTool is the web interface for all communication with the jury. After
login, there will be several buttons to indicate the main parts. Also, in
the upper right corner of the system you'll see a small status frame. The
main parts are:

\begin{citemize}
\item[Home:] a small introduction to the system;

\item[Scorelist:] to see the current contest standings;

\item[Clarifications:] to receive clarifications and send clarification
requests and;

\item[Submissions:] to track the status of your submissions;
\end{citemize}

In the following subsections, we will describe the login process, the status
frame and three of the four main parts. Home is not discussed, as this only
contains a small introduction to TeamTool and no dynamic parts.

\vspace{-0.3cm}
\subsection{Login}
\vspace{-0.2cm}

Your login with your username and password that's provided by the
organisation.

\vspace{-0.3cm}
\subsection{Status frame}
\vspace{-0.2cm}

In the top right corner, the status frame will keep you informed about new
information. It is refreshed automatically and contains four items:

\begin{citemize}
\item The light ``NEW CLARIFICATIONS'' warns you when you receive
clarifications from the jury, if it is yellow you should go to
``Clarifications'' to read them.

\item The light ``NEW RESULTS'' warns you when you have received a new result
for a submission. You can check the submission status by going to
``Submissions''

\item The clock shows the time since the start of the contest. If there is
no clock, but the word ``WAIT'' it means the contest has not started yet. If
the word ``STOP'' appears, the contest has ended. If there's no clock, you
cannot submit files. You can send clarification requests, if you need to.

\item Finally it shows your team rank. Once the scorelist is no longer
updated (usually in the last hour) your rank is no longer updated as well.
The rank considers equal scores and equal time the same rank, so the rank
may be higher than your place in the scorelist. In fact, everyone is first
at the start of the contest.

\end{citemize}

\vspace{-0.3cm}
\subsection{Scorelist}
\vspace{-0.2cm}

The scorelist displays the standings of the contest. It is in general
updated until one hour before the end of the contest and as soon as any
submission is accepted. You might have to refresh this page to see the
actual score.

Standings are decided by:

\begin{citemize}
\item the total number of problems solved (more is better) and;
\item the total time used for the solved problems (less is better).
\end{citemize}

The time is calculated by the time between the contest start and the time of
submission. Every wrong submission gives you 20 penalty minutes, but they
only count if you eventually solve the problem.

\vspace{-0.3cm}
\subsection{Clarifications}
\vspace{-0.2cm}

In the clarifications part, you can send clarification requests or read
clarifications from the jury. A clarification or clarification request can
refer to a specific problem, but this is not required.

The subject can be up to 40 characters, the body doesn't have a low
limit on the number of characters.

On the right of this page is the index of messages. You can read your own
messages sent to the jury, read all clarifications, or just those that refer
to a specific problem. The number of messages is indicated after each group.
The number of unread messages is written between brackets. A message is
marked ``read'' when you click the subject of this message. It then shows
the message instead of the send message window and you can read the message.
It is advised to read clarifications from the jury when they arrive, they
might contain information you need to solve a problem.

If there are unread messages you will be notified by the yellow ``NEW
CLARIFICATIONS'' field in the top right corner. This is refreshed every
thirty seconds.

\vspace{-0.3cm}
\subsection{Submissions}
\vspace{-0.2cm}

The submissions part enables you to submit solutions and see the results of
your submissions. You are advised not to use TeamTool's submission ability,
if possible. \cmd{submit} (see section~\ref{sec:submit}) is the recommended
way, since it has a few more checks to prevent accidentally submitting the
wrong file.

If you do submit using the teamtool, make sure to fill in the right problem
and compiler. Also, \emph{do not refresh the page after you submitted}, this
will result in submitting the file again.

The right side of the window shows the statuses of your own submissions. You
can select to see all submissions or only those of a specified problem. The
status can be:

\begin{citemize}
\item[Pending...:] waiting for a result;
\item[ACCEPTED:] your submission was accepted as being correct;
\item[COMPILER\_ERROR:] the system was not able to compile the file;
\item[RUNTIME\_ERROR:] execution of your submission resulted in a run-time
error;
\item[RUNTIME\_EXCEEDED:] the program did not terminate within the set
timelimit for this problem;
\item[WRONG\_OUTPUT:] the program did not give the correct output or;
\item[NO\_OUTPUT:] the program did not produce any output at all.
\end{citemize}

\strong{Note} You can get RUNTIME\_ERROR in several situations. Your program
may have crashed, it may not have returned a zero exit status or it may have
exceeded limits (such as memory usage or disk quota) set by the system.
\emph{Please make sure your program returns a zero exit status.} Several
languages do this automatically (Java for example), but some
don't (most notably C and C++).

\end{document}
